\section*{AI工具使用声明}
\nocite{deepseek2026,lightgbm2017}

本参赛队在研究过程中使用了以下人工智能工具作为辅助手段。所有AI生成的内容均经作者严格核对、修改和验证，其正确性、可靠性及学术规范性由本参赛队承担全部责任。核心建模、算法设计、代码实现、结果分析与结论由队员独立完成。

具体使用情况如下：

\begin{itemize}[itemsep=4pt]
\item \textbf{DeepSeek-V4-Pro（v4，DeepSeek）}：用于数据 Schema 校验与清洗流水线设计、探索性数据分析（EDA）可视化脚本生成、稀疏基线与消融实验代码编写、严格嵌套级联方案实现，以及论文 LaTeX 结构与部分章节草稿的辅助生成。使用日期：2026年8月30日。

\item \textbf{Opencode（DeepSeek-V4-Pro，v4）}：作为交互式 CLI 编程辅助工具，在代码生成、实验运行与论文撰写环节提供辅助。仅作辅助作用，所有AI生成的内容均经作者严格核对、修改和验证，其正确性、可靠性及学术规范性由本参赛队承担全部责任。核心建模、算法设计、代码实现、结果分析与结论由队员独立完成。使用日期：2026年8月30日。
\end{itemize}

各章节中凡由AI辅助生成并修改后采用的内容，均已在对应位置以脚注标注（注明工具名称、提示词原文、生成日期及用途）。

本参赛队保证：全文AIGC生成比例不超过30\%，论文整体相似比不超过50\%。支撑材料中附有《AI工具使用详情说明》（文件名：BDATA2600409\_AI使用说明.pdf），包含工具基本信息、使用目的、关键交互记录及采纳修改说明。
